% The cost argument. New file, 2026-10-25, for the restarted framing: the
% introduction's "cost prohibitive" is currently an assertion, and this is
% where it becomes arithmetic.
%
% Deliberately a scaffold. The figures, the identity and the numbers are here
% because they are infrastructure; the argument is yours to write. The same
% material is worked through at length in campaigns.pdf \S11 if you want a
% draft to react to rather than a blank page.
\section{What a combinatorial library costs}
\label{sec:res:cost}

\plan{\textbf{Yours to write.} The job of this section, given where it now
sits: Section~\ref{sec:res:ranking} has just established that nothing ranks the
candidates, so the obvious move is to order them all. This section is why that
is not available, and why the combinatorial form is not merely cheaper but
\emph{differently} priced --- exponential in $k$ where a one-sequence-per-well
library is linear.

The one-line version: a million unique full-length genes is \$42M; the same
million as a $k = 6$ combinatorial library is under \$500.}

A protein of $L$ amino acids is $3L$ bases. Cut it into $k$ equal fragments and
each fragment is $3L/k$ bases, so a budget $B$ at a price $c$ per base buys

\begin{equation}
  v \;=\; \frac{B}{3Lc}
  \qquad\text{variants per slot,}\qquad
  S \;=\; v^{k} \;=\; \left(\frac{B}{3Lc}\right)^{k}
  \qquad\text{recombinants.}
  \label{eq:cost}
\end{equation}

\noindent
There is no $k$ in $v$: cutting into $k$ pieces buys $k$ times as many
fragments, each $1/k$ as long, and the two cancel. Taking logs,
$\log S = k(\log B - \log B_1)$ for $B_1 = 3Lc$ the cost of one full-length
gene, so on logarithmic axes the model is a fan of straight lines of slope $k$,
all pivoting on $B_1$.

\begin{figure}[t]
  \centering
  \includegraphics[width=\textwidth]{figures/pdf/panel_cost_size.pdf}
  \caption{Left: Equation~\ref{eq:cost} drawn. $k = 1$ is the line without
  recombination --- buy $v$ genes, have $v$ sequences --- and the vertical
  distance from it to any other line is what recombination buys for the same
  money. Because the axes are logarithmic a different price per base slides the
  pivot sideways and changes nothing else. Right: the same model inverted ---
  what a target library size costs, and how much of that is the vendor's order
  minimum rather than the model.}
  \label{fig:cost-size}
\end{figure}

\begin{figure}[t]
  \centering
  \includegraphics[width=\textwidth]{figures/pdf/panel_cost_k.pdf}
  \caption{Vendor terms, which Equation~\ref{eq:cost} ignores. A fragment is
  billed at $\max(m,\, 3L/k + h)$ for an order minimum $m$ and $h$ bases of
  Golden Gate sites and overhangs. On a \num{200}-aa enzyme the minimum binds
  from $k = 3$ onward and by $k = 8$ three quarters of every base bought is
  padding. The library still grows with $k$ --- the exponent beats the division
  --- so cost does not select a $k$.}
  \label{fig:cost-k}
\end{figure}

\plan{\textbf{Numbers that are settled}, at Twist eBlock pricing
(\$0.07 per base, confirmed) for a \num{200}-aa enzyme:
\begin{planlist}
  \item One full-length gene is \$42. A \$\num{10000} order therefore stocks
        \num{238} variants in every slot, \emph{however finely it is cut}.
  \item At $k = 6$ that is $\num{1.8e14}$ recombinants, against \num{238}
        sequences without recombination.
  \item Reaching $10^{12}$ recombinants costs \$\num{42000} at $k = 4$ and
        \$\num{1329} at $k = 8$.
\end{planlist}

\textbf{And the honest limits}, which belong in the text rather than a
footnote. The order minimum $m = \num{300}$ bases and adapter overhead
$h = \num{40}$ bases are \emph{assumed} and should be checked against a live
quote; the price per base is not assumed. Nothing in this section is measured
--- it is arithmetic on a price --- and it says nothing about whether a library
\emph{works}, which is Section~\ref{sec:res:fragments}. Two libraries of
identical size differ by 50 points of pass rate.

\textbf{What binds before the budget does}, all three worth stating because
they are the reason this section is not the end of the argument: stocking $v$
variants needs $v$ passing designs, so the \num{256}-parent libraries run out
at \$\num{10752}; one-pot assembly gives ${\sim}10^{9}$ transformants; and the
one screen this project has run read \num{3135} wells. Synthesis sets the
ceiling. It has never set the floor --- which is what the rest of the paper is
about.}
